% Build in this directory (run twice for page references):
% pdflatex -interaction=nonstopmode -halt-on-error resume6PT26.tex
% Topics follow the syllabus in ../class1/class1PT26.tex.
% Content adapted from ../lectureNote.tex and ../lectureNotes.tex.
\documentclass[a4paper,10pt,pdflatex,ja=standard,jafont=ipaex-type1,
  magstyle=real,scale=1]{bxjsarticle}
\usepackage{lmodern}
\usepackage{amsmath,amssymb,mathtools}
\usepackage{geometry}
\geometry{reset,left=22mm,right=22mm,top=20mm,bottom=17mm,
  headheight=13pt,headsep=6mm,footskip=10mm}
\usepackage{xcolor}
\usepackage{enumitem}
\usepackage{array,booktabs,tabularx}
\usepackage{tikz}
\usepackage{fancyhdr}
\usepackage{hyperref}
\definecolor{ink}{HTML}{193A5C}
\definecolor{accent}{HTML}{237D82}
\definecolor{soft}{HTML}{EAF3F3}
\definecolor{muted}{HTML}{52616B}
\newcommand{\coursename}{確率論と統計学 I}
\hypersetup{unicode=true,colorlinks=true,urlcolor=accent,linkcolor=ink,
  pdfauthor={中島秀太},pdftitle={第6回：二項分布の正規近似}}
\pagestyle{fancy}
\fancyhf{}
\fancyhead[L]{\small\color{ink}\coursename{}・授業要約と演習}
\fancyhead[R]{\small\color{ink}第6回}
\fancyfoot[C]{\small\color{ink}\thepage/\pageref*{lastpage}}
\renewcommand{\headrulewidth}{0.3pt}
\setlength{\parindent}{0pt}
\setlength{\parskip}{2pt}
\setlength{\abovedisplayskip}{4pt}
\setlength{\belowdisplayskip}{4pt}
\setlength{\abovedisplayshortskip}{3pt}
\setlength{\belowdisplayshortskip}{4pt}
\setlist[itemize]{leftmargin=1.8em,itemsep=2pt,topsep=3pt,parsep=0pt}
\setlist[enumerate]{leftmargin=2em,itemsep=6pt,topsep=4pt,parsep=0pt}
\newcommand{\topic}[1]{\par\vspace{5pt}
  {\large\sffamily\mdseries\color{accent}#1}\par\vspace{2pt}}
\newcommand{\handouttitle}[1]{\begin{center}
  {\Large\color{ink}#1}\par\vspace{3mm}
  {\color{accent}\rule{0.88\linewidth}{0.7pt}}
  \end{center}\vspace{-2mm}}
\newcommand{\keybox}[1]{\par\vspace{3pt}\begingroup
  \setlength{\fboxsep}{5pt}
  \colorbox{soft}{\parbox{\dimexpr\linewidth-2\fboxsep\relax}{#1}}
  \endgroup\par\vspace{2pt}}
\newcommand{\examplelabel}[1]{{\sffamily\mdseries\color{ink}#1}}
\newcommand{\columnheading}[1]{\par\vspace{5pt}
  {\fontsize{11pt}{15pt}\selectfont\sffamily\mdseries\color{accent}#1}\par\vspace{2pt}}
\newcommand{\columnstart}[1]{\clearpage
  \fancyhead[L]{\small\color{ink}\coursename{}・発展的内容}
  {\small\sffamily\mdseries\color{accent}発展的内容}\par\vspace{4pt}
  {\Large\color{ink}#1}\par\vspace{2pt}
  {\color{accent}\rule{\linewidth}{0.7pt}}\par\vspace{4pt}}
\newcommand{\F}{\mathcal F}
\newcommand{\R}{\mathbb R}
\newcommand{\N}{\mathbb N}
\newcommand{\E}{\mathbb E}
\newcommand{\Var}{\operatorname{Var}}
\newcommand{\Cov}{\operatorname{Cov}}
\newcommand{\ind}{\mathrm{1}}
\newcommand{\iid}{\text{独立同分布}}
\begin{document}
\fontsize{10pt}{13.5pt}\selectfont
\handouttitle{第6回\quad 二項分布の正規近似}
\keybox{二項分布の中心は $n/2$、幅は $\sqrt n/2$ である。
中心と幅をそろえると、離散的な確率質量の和が正規密度の積分へ近づく。}

\topic{公平なコイントスと標準化}
$X_1,X_2,\ldots$ は独立で $P(X_j=0)=P(X_j=1)=1/2$ とする。
表の個数 $S_n=\sum_{j=1}^nX_j$ について
$$P(S_n=k)=\binom nk2^{-n},\qquad
 \E[S_n]=\frac n2,\qquad \Var(S_n)=\frac n4.$$
標準化した変数と格子幅は
$$Z_n:=\frac{S_n-n/2}{\sqrt n/2},\qquad
 x_{n,k}:=\frac{2k-n}{\sqrt n},\qquad \Delta x=\frac2{\sqrt n}.$$
$Z_n$ の平均は $0$、分散は $1$ となる。
値の刻みは $n$ とともに小さくなるが、各 $n$ で $Z_n$ は離散型である。

\topic{正規密度と Gauss 積分}
標準正規密度とその分布関数を
$$\phi(x):=\frac1{\sqrt{2\pi}}e^{-x^2/2},\qquad
 \Phi(t):=\int_{-\infty}^t\phi(x)\,dx$$
とおく。正規化のために $I=\int_\R e^{-x^2/2}\,dx$ を計算する。
非負積分に Tonelli の定理を使い、極座標に変換すれば
$$I^2=\int_{\R^2}e^{-(x^2+y^2)/2}\,dx\,dy
 =2\pi\int_0^\infty r e^{-r^2/2}\,dr=2\pi.$$
$I>0$ なので $I=\sqrt{2\pi}$、したがって $\int_\R\phi=1$ となる。

\topic{Stirling 公式から中心の高さを求める}
階乗の漸近公式
$$n!\sim\sqrt{2\pi n}\left(\frac ne\right)^n$$
を使う。偶数 $n$ について中心の確率は
$$P(S_n=n/2)=\frac{n!}{((n/2)!)^2}2^{-n}
 \sim\sqrt{\frac2{\pi n}}=\phi(0)\frac2{\sqrt n}.$$
中心の一点の確率は $n^{-1/2}$ 程度となる。
$n=100$ なら厳密値は約 $0.079589$、上の近似値は約 $0.079788$ である。
密度の高さそのものに比べるのではなく、格子幅を掛けて比べる点が重要である。

\topic{確率質量と正規密度の比較（$n=20$）}
\begin{center}
\begin{tikzpicture}[x=0.43cm,y=14cm,font=\small]
 \foreach \k/\p in {0/0.000000954,1/0.000019073,2/0.000181198,
 3/0.001087189,4/0.004620552,5/0.014785767,6/0.036964417,
 7/0.073928833,8/0.120134354,9/0.160179138,10/0.176197052,
 11/0.160179138,12/0.120134354,13/0.073928833,14/0.036964417,
 15/0.014785767,16/0.004620552,17/0.001087189,18/0.000181198,
 19/0.000019073,20/0.000000954}{
   \fill[ink!22] ({\k-0.45},0) rectangle ({\k+0.45},\p);
 }
 \draw[accent,line width=0.9pt,domain=0:20,samples=101]
  plot (\x,{exp(-(\x-10)^2/10)/sqrt(10*pi)});
 \draw[ink,->] (-0.7,0)--(21,0) node[right] {$k$};
 \draw[ink,->] (0,0)--(0,0.22);
 \foreach \k in {0,5,10,15,20}{\node[below,ink] at (\k,0) {$\k$};}
 \foreach \p in {0.1,0.2}{
   \draw[muted!35,dotted] (0,\p)--(20.5,\p);
   \node[left,ink] at (0,\p) {$\p$};
 }
 \node[accent,right] at (12.3,0.19) {$N(10,5)$ の密度};
\end{tikzpicture}
\end{center}
{\small\color{muted}棒の高さは $P(S_{20}=k)$。曲線は平均10・分散5の正規密度で、
一点の確率は幅1の区間の面積で近似する。}\par

\newpage
\topic{中心付近の局所近似}
固定した $R>0$ に対し、$|x_{n,k}|\le R$ の範囲では
$$P(S_n=k)=\phi(x_{n,k})\frac2{\sqrt n}\,(1+o(1))$$
が $k$ について一様に成り立つ。
Stirling 公式を二項係数に代入し、$u=(2k-n)/n=x_{n,k}/\sqrt n$ とおくと、
指数部分には
$$\frac{(1+u)\log(1+u)+(1-u)\log(1-u)}2
 =\frac{u^2}2+O(u^4)$$
が現れる。$x_{n,k}$ が有界なら $nu^2=x_{n,k}^2$、$nu^4\to0$ なので
$e^{-x_{n,k}^2/2}$ が残る。係数は $\sqrt{2/(\pi n)}$ に近づく。

\topic{確率質量を足すと積分になる}
固定した有限区間 $[a,b]$ では、上の一様近似と Riemann 和により
$$P(a\le Z_n\le b)
 =\sum_{a\le x_{n,k}\le b}P(S_n=k)
 \longrightarrow\int_a^b\phi(x)\,dx.$$
無限区間に進むには尾を抑える。$\E[Z_n^2]=1$ なので
$$P(|Z_n|>R)\le\frac1{R^2}.$$
まず有限区間で $n\to\infty$、次に $R\to\infty$ として
$$P(Z_n\le t)\longrightarrow\Phi(t)\qquad(t\in\R)$$
を得る。これが公平なコイントスに対する de Moivre--Laplace の定理である。

\topic{実際の計算：連続性補正}
整数 $a\le b$ に対し、$S_n$ の一点 $k$ を区間 $[k-1/2,k+1/2]$ に対応させると
$$P(a\le S_n\le b)\approx
 \Phi\left(\frac{b+1/2-n/2}{\sqrt n/2}\right)
 -\Phi\left(\frac{a-1/2-n/2}{\sqrt n/2}\right).$$
たとえば $n=100$、$a=45$、$b=55$ では $\Phi(1.1)-\Phi(-1.1)$ となる。
これは有限 $n$ での精度を改善する方法であり、極限定理の主張そのものとは区別する。

\topic{一般の成功確率の場合}
固定した $0<p<1$ に対しても、独立な $\mathrm{Bern}(p)$ の和は
$$\frac{S_n-np}{\sqrt{np(1-p)}}\ \Rightarrow\ N(0,1).$$
$\Rightarrow$ は分布収束を表す。一般の独立同分布列への拡張は第10回に扱う。
小さい $n$ や極端な尾では、正規近似の誤差に注意する。

\newpage
\topic{演習問題}
\begin{enumerate}[label=\arabic*.]
\item 独立な $X_j\sim\mathrm{Bern}(p)$ に対し、$\E[S_n]$ と $\Var(S_n)$ を求めよ。
標準化で平均 $0$、分散 $1$ になることを確かめよ。
\item $S_{100}\sim\mathrm{Bin}(100,1/2)$ について、$P(45\le S_{100}\le55)$ を
連続性補正あり・なしで正規近似せよ。厳密な二項和と比べよ。
\item $Z_n$ の局所近似を用いて、有限な $a<b$ に対する区間確率の収束を示せ。
各項の近似誤差を足し合わせるため、一様性をどこで使うか説明せよ。
\item $P(|Z_n|>R)\le R^{-2}$ を証明し、有限区間での収束から分布関数の収束へ
進む際の誤差を評価せよ。
\end{enumerate}

\topic{レポート問題}
\begin{enumerate}[label=\arabic*.]
\item $n=100$ の公平なコイントスを $10\,000$ 回反復するシミュレーションを行え。
$S_n$ の確率質量のヒストグラムに、$N(50,25)$ の密度を重ねよ。
標準化して描く場合は、密度と確率質量のどちらを表示しているか明記せよ。
\item $n=20,50,200$ の結果を比較し、中心部と尾部の近似誤差、
および連続性補正の効果を考察せよ。試行回数を増やしても残る理論上の近似誤差と、
反復回数による標本の揺らぎを区別して述べよ。
\end{enumerate}

\topic{考えるための手がかり}
幅 $h$ のビンの度数を全反復回数だけで割ると確率質量、さらに $h$ で割ると密度の推定になる。
$Z_n$ の格子幅は $2/\sqrt n$ なので、正規密度と確率質量を直接重ねると縦軸が合わない。
数値実験を行う場合は、使用した乱数生成方法と反復回数も記録しよう。
\vfill
{\small\color{muted}解答メモ：中心化、標準化、格子幅、端点の補正を順に確認する。}

\columnstart{コインの足し算から、ランダムな軌跡へ}
コインの表の個数を数える代わりに、表なら右へ一歩、裏なら左へ一歩進む。
同じコイントスの列が、確率分布だけでなく時間とともに動く道のモデルを与える。

\columnheading{位置の大きさは $\sqrt n$ 程度}
$Y_j=2X_j-1\in\{-1,1\}$、$R_m=\sum_{j=1}^mY_j$ とおく。
独立性と対称性から
$$\E[R_m]=0,\qquad \Var(R_m)=m,\qquad
 \frac{R_n}{\sqrt n}=Z_n\Rightarrow N(0,1).$$
一歩ずつは同じ長さでも、右と左が打ち消し合うため、典型的な位置の幅は $\sqrt n$ となる。
一方向へ進み続ける場合の距離 $n$ とは異なる尺度である。

\columnheading{時間と空間を同時に縮める}
時間間隔を $1/n$、歩幅を $1/\sqrt n$ とし、格子時刻で
$$W^{(n)}(k/n)=\frac{R_k}{\sqrt n}$$
と定め、格子点の間を線形に補間する。
各時刻 $t$ での分散はほぼ $t$、増分の分散はほぼ経過時間となる。
たとえば補間前の値について $s\le t$ なら
$$\Cov\left(\frac{R_{\lfloor ns\rfloor}}{\sqrt n},
 \frac{R_{\lfloor nt\rfloor}}{\sqrt n}\right)
 =\frac{\lfloor ns\rfloor}{n}\longrightarrow s=\min\{s,t\}.$$
過去までに共有した歩数が、二つの時刻の位置の相関を決めている。

\columnheading{極限としての Brown 運動}
Brown 運動 $(B_t)_{t\ge0}$ は $B_0=0$、連続な軌跡、独立な増分をもち、
$B_t-B_s\sim N(0,t-s)$（$0\le s<t$）を満たす過程である。
一時刻の中心極限定理を、連続な軌跡全体の収束へ強めると、
補間した $W^{(n)}$ は $C([0,1])$ 上で Brown 運動へ弱収束する。
この主張は Donsker の不変原理の、公平な符号列に対する場合である。

\columnheading{一時刻の定理と軌跡の定理の違い}
各固定時刻の分布が近づくだけでは、軌跡全体が近づくとは限らない。
複数時刻の同時分布に加え、短い時間で大きく振動しすぎないことも制御する必要がある。
本講義の中心極限定理はその入口である。
分布の収束を学ぶときには、何がランダムな対象なのかが、数から関数へ広がる点も意識してみよう。

\label{lastpage}
\end{document}
